\subsection{连接同态的计算}

考虑一个交换图

\begin{enumerate}
\def\labelenumi{(\arabic{enumi})}
\tightlist
\item
\end{enumerate}

\[
\begin{array}{c} 0 \longrightarrow \mathrm{R} ^ {\prime} \xrightarrow {\tilde {u}} \mathrm{R} \xrightarrow {\tilde {v}} \mathrm{R} ^ {\prime \prime} \longrightarrow 0 \\ a ^ {\prime} \Big \downarrow \qquad \qquad \qquad \qquad \qquad \qquad \qquad \qquad \qquad \qquad \qquad \qquad \qquad \qquad \qquad \qquad \qquad \qquad \qquad \qquad \qquad \qquad \qquad \qquad \qquad \qquad \qquad \qquad \qquad \qquad \qquad \qquad \qquad \qquad \\ 0 \longrightarrow \mathrm{M} ^ {\prime} \xrightarrow {u} \mathrm{M} \xrightarrow {v} \mathrm{M} ^ {\prime \prime} \longrightarrow 0 \end{array}\tag{2}
\]

其中第一行 (1) 是左 A-模复形的正合列，第二行 (2) 是左 A-模的正合列，且竖直箭头为左分解。

命题 3. --- a) 设 P 是一个 A-模， \(b: S \to P\) 为 P 的一个左分解；假设 k-模复形序列

\[
0 \rightarrow \mathrm{S} \otimes_ {\mathrm{A}} \mathrm{R} ^ {\prime} \xrightarrow {l _ {\mathrm{s}} \otimes \tilde {u}} \mathrm{S} \otimes_ {\mathrm{A}} \mathrm{R} \xrightarrow {l _ {\mathrm{s}} \otimes \tilde {v}} \mathrm{S} \otimes_ {\mathrm{A}} \mathrm{R} ^ {\prime \prime} \rightarrow 0\tag{3}
\]

是正合的。则下面的图表是交换的：

\[
\begin{array}{c c c} \operatorname{Tor} ^ {\mathbf {A}} (\mathrm{P}, \mathrm{M} ^ {\prime \prime}) & \xrightarrow {\partial (\mathrm{P} , (2))} & \operatorname{Tor} ^ {\mathbf {A}} (\mathrm{P}, \mathrm{M} ^ {\prime}) \\ \psi (\mathrm{S}, \mathrm{R} ^ {\prime \prime}) \Big \downarrow & & \Big \downarrow \psi (\mathrm{S}, \mathrm{R} ^ {\prime}) \\ \mathrm{H} (\mathrm{S} \otimes_ {\mathbf {A}} \mathrm{R} ^ {\prime \prime}) & \xrightarrow {\partial ((3))} & \mathrm{H} (\mathrm{S} \otimes_ {\mathbf {A}} \mathrm{R} ^ {\prime}). \end{array}
\]

\begin{enumerate}
\def\labelenumi{\alph{enumi})}
\setcounter{enumi}{1}
\tightlist
\item
  设 N 是一个左 A-模， \(c : N \to E\) 为 N 的一个右分解；假设 k-模复形序列
\end{enumerate}

\begin{enumerate}
\def\labelenumi{(\arabic{enumi})}
\setcounter{enumi}{3}
\tightlist
\item
\end{enumerate}

\(0 \to \text{Homgr}_{\mathbf{A}}(\mathbf{R}^{\prime \prime}, \mathbf{E}) \xrightarrow{\text{Homgr}_{\mathbf{A}}(\tilde{\imath}, 1)} \text{Homgr}_{\mathbf{A}}(\mathbf{R}, \mathbf{E}) \xrightarrow{\text{Homgr}_{\mathbf{A}}(\tilde{\mu}, 1)} \text{Homgr}_{\mathbf{A}}(\mathbf{R}^{\prime}, \mathbf{E}) \to 0\) 是正合的。则下面的图表是交换的：

\[
\begin{array}{c} \mathrm{H} (\operatorname{Homgr} _ {\mathbf {A}} (\mathrm{R} ^ {\prime}, \mathrm{E})) \xrightarrow {\delta ((4))} \mathrm{H} (\operatorname{Homgr} _ {\mathbf {A}} (\mathrm{R} ^ {\prime \prime}, \mathrm{E})) \\ \varphi (\mathrm{R} ^ {\prime}, \mathrm{E}) \Big \downarrow \\ \operatorname{Ext} _ {\mathbf {A}} (\mathrm{M} ^ {\prime}, \mathrm{N}) \xrightarrow {\partial ((2) , \mathrm{N})} \operatorname{Ext} _ {\mathbf {A}} (\mathrm{M} ^ {\prime \prime}, \mathrm{N}). \end{array}
\]

例如证明 \(a\) )。设 \(\beta: L(P) \to S\) 是复形的态射，使得 \(b \circ \beta = p_{P}\) ；考虑 \(k\) -复形的图表

\[
\begin{array}{c}0 \rightarrow \mathrm{S} \otimes_ {\mathrm{A}} \mathrm{R} ^ {\prime} \xrightarrow {1 \otimes \tilde {u}} \mathrm{S} \otimes_ {\mathrm{A}} \mathrm{R} \xrightarrow {1 \otimes \tilde {v}} \mathrm{S} \otimes_ {\mathrm{A}} \mathrm{R} ^ {\prime \prime} \longrightarrow 0\\\beta \otimes 1 _ {\mathrm{R}} \Bigg | \quad \beta \otimes 1 _ {\mathrm{R}} \Bigg | \quad \beta \otimes 1 _ {\mathrm{R}} \Bigg |\\0 \rightarrow \mathrm{L(P)} \otimes_ {\mathrm{A}} \mathrm{R} ^ {\prime} \xrightarrow {1 \otimes \tilde {u}} \mathrm{L(P)} \otimes_ {\mathrm{A}} \mathrm{R} \xrightarrow {1 \otimes \tilde {v}} \mathrm{L(P)} \otimes_ {\mathrm{A}} \mathrm{R} ^ {\prime \prime} \longrightarrow 0\\1 \otimes a ^ {\prime} \Bigg | \quad 1 \otimes a \Bigg | \quad 1 \otimes a ^ {\prime \prime} \Bigg |\\0 \rightarrow \mathrm{L(P)} \otimes_ {\mathrm{A}} \mathrm{M} ^ {\prime} \xrightarrow {1 \otimes u} \mathrm{L(P)} \otimes_ {\mathrm{A}} \mathrm{M} \xrightarrow {1 \otimes v} \mathrm{L(P)} \otimes_ {\mathrm{A}} \mathrm{M} ^ {\prime \prime} \longrightarrow 0.\end{array}
\]

它是交换的且各行正合（第一行由假设，另外两行由 L(P) 是平坦的这一事实）。因此我们有交换图（X，第 31 页，命题 2，以及 X，第 72 页，定义 2）。

\[
\begin{array}{c c c} \mathrm{H} (\mathrm{S} \otimes_ {\mathrm{A}} \mathrm{R} ^ {\prime \prime}) & \xrightarrow {\partial ((3))} & \mathrm{H} (\mathrm{S} \otimes_ {\mathrm{A}} \mathrm{R} ^ {\prime}) \\ \mathrm{H} (\beta \otimes 1) \uparrow & & \uparrow \mathrm{H} (\beta \otimes 1) \\ \mathrm{H} (\mathrm{L} (\mathrm{P}) \otimes_ {\mathrm{A}} \mathrm{R} ^ {\prime \prime}) & \longrightarrow & \mathrm{H} (\mathrm{L} (\mathrm{P}) \otimes_ {\mathrm{A}} \mathrm{R} ^ {\prime}) \\ \mathrm{H} (1 \otimes a ^ {\prime \prime}) \downarrow & & \downarrow \mathrm{H} (1 \otimes a ^ {\prime}) \\ \mathrm{H} (\mathrm{L} (\mathrm{P}) \otimes_ {\mathrm{A}} \mathrm{M} ^ {\prime \prime}) & \longrightarrow & \mathrm{H} (\mathrm{L} (\mathrm{P}) \otimes_ {\mathrm{A}} \mathrm{M} ^ {\prime}) \\ \psi_ {\mathrm{P}} (\mathbf {M} ^ {\prime \prime}) \uparrow & & \uparrow \psi_ {\mathrm{P}} (\mathbf {M} ^ {\prime}) \\ \operatorname{Tor} (\mathrm{P}, \mathbf {M} ^ {\prime \prime}) & \xrightarrow {\partial (\mathrm{P} , (2))} & \operatorname{Tor} (\mathrm{P}, \mathbf {M} ^ {\prime}). \end{array}
\]

由 X，第 67 页，命题 4，H(1 ⊗ a′′) 和 H(1 ⊗ a′) 是双射；另一方面，根据同态 ψ 的定义，我们有 H(β ⊗ 1) ∘ ψ(L(P), R′′) = ψ(S, R′′) 以及 H(1 ⊗ a′′)∘ψ(L(P), R′′) = ψ(L(P), M′′) = ψ\_P(M′′)，因此

\[
\psi (\mathrm{S}, \mathrm{R} ^ {\prime \prime}) = \mathrm{H} (\beta \otimes 1) \circ \mathrm{H} (1 \otimes a ^ {\prime \prime}) ^ {- 1} \circ \psi_ {\mathrm{P}} (\mathrm{M} ^ {\prime \prime});
\]

类似地， \(\psi(S, R') = H(\beta \otimes 1) \circ H(1 \otimes a')^{-1} \circ \psi_P(M')\) ，而所需的结论 \(\partial((3)) \circ \psi(S, R'') = \psi(S, R') \circ \partial(P, (2))\) 由上面图表的交换性得出。

注记. --- 1) 利用交换同构，由 a) 可推出将张量积的两个变元互换角色后所得的类似命题。

\begin{enumerate}
\def\labelenumi{\arabic{enumi})}
\setcounter{enumi}{1}
\tightlist
\item
  沿用 \(a\) ) 的记号，假设 S 平坦，或 R、R'、R'' 平坦；则一方面序列 (3) 是正合的（X，第 72 页，推论 2），从而可应用命题 3；另一方面 \(\psi(S, R')\) 是双射的（定理 1），因此
\end{enumerate}

\[
\partial (\mathrm{P}, (2)) = \psi (\mathrm{S}, \mathrm{R} ^ {\prime}) ^ {- 1} \circ \partial ((3)) \circ \psi (\mathrm{S}, \mathrm{R} ^ {\prime \prime}).
\]

\begin{enumerate}
\def\labelenumi{\arabic{enumi})}
\setcounter{enumi}{2}
\tightlist
\item
  沿用 \(b\) ) 的记号，假设 E 是内射的，或者 R、R'、R'\,' 是投射的；则一方面序列 (4) 是正合的（X，第 83 页，命题 2），可以应用命题 3；另一方面， \(\varphi(R', E)\) 是双射的（定理 1）；因此
\end{enumerate}

\[
\delta ((2), \mathrm{N}) = \varphi (\mathrm{R} ^ {\prime \prime}, \mathrm{E}) \circ \hat {\partial} ((4)) \circ \varphi (\mathrm{R} ^ {\prime}, \mathrm{E}) ^ {- 1}.
\]

现在考虑一个交换图，其第一行 (5) 是左 A-模的正合列，第二行 (6) 是左 A-模复形的正合列，且竖直箭头是右分解。如同命题 3 中的推理，可证明以下命题：

\begin{enumerate}
\def\labelenumi{(\arabic{enumi})}
\setcounter{enumi}{4}
\tightlist
\item
\end{enumerate}

\[
\begin{array}{c} 0 \longrightarrow \mathrm{N} ^ {\prime} \xrightarrow {r} \mathrm{N} \xrightarrow {s} \mathrm{N} ^ {\prime \prime} \longrightarrow 0 \\ c ^ {\prime} \Big \downarrow \qquad \qquad c \Big \downarrow \qquad \qquad c ^ {\prime \prime} \Big \downarrow \\ 0 \longrightarrow \mathrm{E} ^ {\prime} \xrightarrow {\tilde {r}} \mathrm{E} \xrightarrow {\tilde {s}} \mathrm{E} ^ {\prime \prime} \longrightarrow 0 \end{array}\tag{6}
\]

命题 4. --- 设 M 是左 A-模，

且 \(a: \mathbf{R} \to \mathbf{M}\) 为 M 的一个左分解，使得序列

\begin{enumerate}
\def\labelenumi{(\arabic{enumi})}
\setcounter{enumi}{6}
\tightlist
\item
  \(0 \to \operatorname{Homgr}_{\mathbf{A}}(\mathbf{R}, \mathbf{E}') \xrightarrow{\operatorname{Homgr}(\tilde{\mathcal{S}}, 1)} \operatorname{Homgr}_{\mathbf{A}}(\mathbf{R}, \mathbf{E}) \xrightarrow{\operatorname{Homgr}(\tilde{\mathcal{F}}, 1)} \operatorname{Homgr}_{\mathbf{A}}(\mathbf{R}, \mathbf{E}'' \to 0\) 是正合的。则下图是交换的。
\end{enumerate}

\[
\begin{array}{c} \mathrm{H} (\operatorname{Homgr} _ {\mathbf {A}} (\mathbf {R}, \mathbf {E} ^ {\prime \prime})) \xrightarrow {\partial ((7))} \mathrm{H} (\operatorname{Homgr} _ {\mathbf {A}} (\mathbf {R}, \mathbf {E} ^ {\prime})) \\ \varphi (\mathbf {R}, \mathbf {E} ^ {\prime \prime}) \Biggl \downarrow \\ \operatorname{Ext} _ {\mathbf {A}} (\mathbf {M}, \mathbf {N} ^ {\prime \prime}) \xrightarrow {\delta (\mathbf {M} , (5))} \operatorname{Ext} _ {\mathbf {A}} (\mathbf {M}, \mathbf {N} ^ {\prime}). \end{array}
\]

注记. --- 4) 若 R 是投射的，或 E、E'、E'' 是内射的，则序列 (7) 是正合的（X，第 83 页，命题 2）；此外，φ(R, E'') 是双射的（定理 1）；因此

\[
\delta (M, (5)) = \varphi (R, E ^ {\prime}) \circ \partial ((7)) \circ \varphi (R, E ^ {\prime \prime}) ^ {- 1}.
\]

\begin{enumerate}
\def\labelenumi{\arabic{enumi})}
\setcounter{enumi}{4}
\tightlist
\item
  我们把与命题 3 和命题 4 类似且关于同态 \(\psi_{1}\) 与 \(\varphi_{1}\) 的命题的陈述与证明留给读者。
\end{enumerate}

